\section{现场速查}

这一章不追求从头到尾读，而是放在手边，遇到问题直接翻到对应条目。

\subsection{现象$\rightarrow$参数对照表}

\begin{table}[H]
    \centering
    \caption{现场现象与参数调节对照}
    \label{tab:symptom-param}
    \begin{tabular}{p{5.5cm}lp{3cm}}
        \toprule
        \textbf{现场现象} & \textbf{优先调整的参数} & \textbf{调节方向} \\
        \midrule
        传送带启动太早，和前工序冲突 & 等待时间 & 增大 \\
        传送带启动太晚，前后工序有空档 & 等待时间 & 减小 \\
        纸盒没送到指定位置 & 运行时间 & 增大 \\
        运行时间太长，影响下一周期 & 运行时间 & 减小 \\
        整体节拍太慢，前工序等不及 & 速度 & 增大 \\
        送料过快/不稳 & 速度 & 减小 \\
        启动瞬间冲击明显（振动/噪音大） & 加速度 & 减小 \\
        启动太慢、拖太久 & 加速度 & 增大 \\
        \bottomrule
    \end{tabular}
\end{table}

\textbf{原则：一次只改一个参数，改完先做"单次运行测试"确认效果，再改下一个。}

\subsection{设备状态说明}

界面右上角显示当前设备模式，不同模式的含义和处理方式：

\begin{table}[H]
    \centering
    \caption{设备状态说明}
    \label{tab:device-status}
    \begin{tabular}{lp{4.5cm}p{5cm}}
        \toprule
        \textbf{显示} & \textbf{含义} & \textbf{售后处理方式} \\
        \midrule
        IDLE & 空闲，可以调参和测试 & 正常操作 \\
        RUNNING & 电机正在执行动作 & 等动作结束后再操作 \\
        SERVICE & 售后临时调试状态 & 正常，程序会自动进入或退出 \\
        ERROR & 异常状态 & 断电重启后重新连接；仍异常则反馈研发 \\
        \bottomrule
    \end{tabular}
\end{table}

售后不需要手动切换模式。点"应用"速度或 PWM 时，程序会自动进入 SERVICE 模式；保存参数后会自动回到 IDLE。

\begin{figure}[H]
    \centering
    \includegraphics[width=0.8\textwidth]{figures/ui-status.png}
    \caption{界面顶部状态栏。显示连接状态和当前设备模式。}
    \label{fig:ui-status}
\end{figure}

\subsection{应用 vs 保存}

\begin{table}[H]
    \centering
    \caption{操作功能对比}
    \label{tab:action-compare}
    \begin{tabular}{lp{4cm}cc}
        \toprule
        \textbf{操作} & \textbf{作用} & \textbf{是否写入设备} & \textbf{断电后是否保留} \\
        \midrule
        保存 & 把参数写进设备 & 是 & 是 \\
        应用（仅限速度、PWM） & 临时生效试一下 & 否 & 否 \\
        读取参数 & 从设备读当前值到界面 & 否 & — \\
        单次运行测试 & 触发一次动作 & 否 & — \\
        \bottomrule
    \end{tabular}
\end{table}

简单记法：
\begin{itemize}
    \item 点"保存"的参数，会写入设备，断电后保留。
    \item 点"应用"的速度或 PWM，主要用于临时测试。只能在 IDLE 状态下使用，会自动进入 SERVICE 模式。
    \item "读取参数"只读取，不修改设备。
    \item "单次运行测试"只触发一次动作，不修改参数。
\end{itemize}

\begin{figure}[H]
    \centering
    \includegraphics[width=0.5\textwidth]{figures/ui-toast.png}
    \caption{保存成功后的提示弹窗，约 1 秒后自动消失。}
    \label{fig:ui-toast}
\end{figure}

\subsection{恢复工程默认值}

如果现场参数被调乱，可以按下面值手动恢复：

\begin{table}[H]
    \centering
    \caption{参数恢复默认值}
    \label{tab:default-values}
    \begin{tabular}{lc}
        \toprule
        \textbf{参数} & \textbf{恢复值} \\
        \midrule
        等待时间 & 250 ms \\
        运行时间 & 250 ms \\
        加速度 & 100 \\
        速度 & 48226 cm/min \\
        电机 A & 1 \\
        电机 B & 0 \\
        \bottomrule
    \end{tabular}
\end{table}

恢复方法：在对应输入框填入上表数值，分别点击对应参数的"保存"按钮。

如果高级调试里的 PWM 被改乱，可把 PWM 改回 \texttt{150} 并点击"保存"。
